\section{Discussion}
\label{sec:discussion}

\subsection{Principal findings}

Separating a rule-based, stored selection from the steps where a model helps made \tool{}'s study selection exactly
reproducible from its stored records, while the model's contribution stayed measurable and bounded. Repeating a run
does not give the same records: the model's search plan varies in its synonyms, and databases differ in what they
return on a given day. The corpus turns this from an unknown into a record: it holds what was found, what was
asked, and every decision, so a review can be audited and its selection replayed even when a new search would
differ. The model's screening was highly consistent with itself on the same records ($\kappa$ above 0.9), and the
score's constants could be changed by half without moving recall by more than a third of a percentage point.

On held-out Cochrane intervention reviews, the systematic-style mode found about half of each review's included
studies and selected a quarter. This is far from what a systematic search achieves, and we do not suggest it can
replace one. The loss analysis shows why, and where to improve: four in five lost studies were lost at search, half
of them because they ranked beyond the 500 PubMed records retrieved by queries that do match them. Ranking and
screening, the steps most LLM tools focus on, accounted for less of the loss; the model screen kept 87\% of the
included studies it saw.

\subsection{What the design costs}

Keeping the selection independent of the model has a price, which the evaluation makes visible. Selection reads a
rule-based parse of the question, and where the parse was poor (a verb or participle joined to a concept, a class
named where studies name its members) included studies were judged off topic even though the model's search plan
had named the concept correctly. Two of 22 reviews selected nothing for this reason. In the evaluation each question
was the review's title, run with the default protocol and no person confirming it; in \tool{}'s guided review the
user sees and edits the parsed concepts before anything is searched, which is the intended remedy. The recall
reported here is therefore that of the tool used without that step.

\subsection{Comparison with other work}

Tools that rank records for screening are usually evaluated on a fixed set of records, by the work saved before a
target recall is reached \citep{omara2015using,vandeschoot2021open}; that measures ranking and not search. Studies of
LLM screening report agreement with human decisions on given records \citep{guo2024automated,khraisha2024can}, and
studies of LLM-written search strings report precision and recall of the strings \citep{wang2023can}. Our evaluation
follows the whole path from question to selected studies against the included studies of published reviews, and
attributes each lost study to a stage, which is what a reviewer needs to judge where the tool can be trusted. The
repeatability problem reported for a commercial LLM assistant \citep{elicit2025repeat} is addressed here not by
making the search deterministic, which the evidence suggests is not possible with live databases and models, but by
storing what each run found and making the selection a replayable function of it.

\subsection{Implications}

For tool developers, the largest gains lie in retrieval depth and query sensitivity rather than ranking: deeper
paging of the queries that already match, and better handling of intervention classes (for example through
pharmacological-action headings). For users, the quick-answer mode of such tools, which held 4\% of a review's
included studies, should not be mistaken for a literature search, and a systematic-style result should be read with
its pool recall in mind. Consistent with recent guidance on AI in evidence synthesis \citep{raise2026}, \tool{}
labels its results, shows its search strings and counts, keeps every screening decision open to the user's override
with a recorded reason, and marks a review that synthesised fewer than all its eligible studies as not complete.

\subsection{Limitations}
\label{sec:limits}

The benchmark covers Cochrane intervention reviews published before mid-2019, and recall may differ for other
questions, fields and review types. Cochrane reviews search sources \tool{} does not (Embase, CENTRAL, trial
registers, grey literature) and screen full texts, so their included studies are a demanding reference standard for
a tool limited to open databases. One model was used throughout; a second screening model could not be measured,
and the model screened titles and abstracts truncated to 420 characters. Search and selection were developed on the
development reviews and frozen after the systematic-style test results had been seen (no change was made to them on
account of those results). Repeats of the record search were made with rule-based planning only. Risk-of-bias
suggestions were not compared with published judgements, because the benchmark does not include them and few
included studies have open full text; they are reported as exploratory. \PENDING{claim-accuracy limitation: single
main annotator, abstracts only, sample of 200 of 881 sentences}. Finally, the software and this evaluation were built
by a single author with extensive AI assistance (Section~\ref{sec:ai}); every result is produced by scripts from
stored runs, and the code, corpora and results are public so that they can be checked.

\subsection{Future work}

Planned work includes deeper and more sensitive searches, a second screening model, a comparison of the risk-of-bias
suggestions with published judgements, and a study of the guided review with reviewers. A count-based evidence gap
map, which fills a grid of interventions by outcomes with dated database counts rather than with the selected
studies \citep{snilstveit2016evidence}, has been added to \tool{} since the evaluation and will be evaluated
separately against published evidence gap maps.

\section{Conclusion}

Making study selection a replayable function of stored records, with a bounded and recorded role for the model,
makes LLM-assisted study selection auditable and reproducible without requiring live searches or models to behave
deterministically. Measured honestly against Cochrane reviews, such a tool finds about half of a review's included
studies, and most of what it misses is lost at search, which points to where these tools should improve and confirms
that they support, rather than replace, a systematic search.
